\section{还没做完的，以及接下来做什么}

\subsection{明确没做的}

以下三项不是"还差一点"，是本次决定不做，理由如下。

**陪练角色的情绪实现。** 设计文档里有完整方案（情绪状态、语气档位、触发条件、禁止句式），代码里没有。题目允许非主角色只做设计，我把时间投给了军师和老师。需要强调的是：如果问"陪练有没有情绪"，准确答案是设计上有、实现上没有。

**联网 PVP。** 本地对战是同一台设备上分屏轮流操作，服务端仍然接受客户端传来的状态快照，不是权威对局服务。真实竞技要求服务端下发可见状态和模式，否则改一个字段就能绕过限制。这条差距不是工程量问题，是架构问题。

**LLM 权重训练。** 本机是 arm64，没有可用的 CUDA 后端，也没有云上 GPU 预算。仓库里记录的 \code{cloud\_budget} 是 \code{not supplied}。

\subsection{受条件限制没做的}

需要外部条件才能完成的，与上面区分开。

\begin{table}[h]
\centering
\small
\begin{tabular}{L{4.2cm}L{9.3cm}}
\toprule
项目 & 缺什么 \\
\midrule
独立大样本模型质量评测 & 独立标注集与 API 预算；现有 6 条真实样本不构成质量结论 \\
检索显著性 & 需要把查询集扩到百条量级，否则差异无法与噪声区分 \\
无提示学习迁移 & 需要真人被试；机制已实现（独立成功与提示后成功分开记录），缺的是数据 \\
慢模型与长局浏览器验收 & 需要可控的慢网环境与完整走查 \\
语音听感 & 需要真人耳朵；本机无音频输出，我无法自验 \\
两名真人多局对战 & 需要两个人；分屏逻辑我用脚本模拟点击验证过 \\
\bottomrule
\end{tabular}
\end{table}

\subsection{接下来的动作}

按可核查的方式写，每项都带完成标志。

\begin{table}[h]
\centering
\small
\begin{tabular}{L{6.2cm}L{4.2cm}L{2.2cm}}
\toprule
动作 & 完成标志 & 预计 \\
\midrule
把评测查询集从 16 条扩到 100 条以上 & 能对关键词与混合方案给出显著性检验结果 & 半天 \\
补一轮 30--50 条真实模型用例 & 产出工具选择正确率与多余调用率 & 半天 \\
给输出校验补因果语义检查 & 已复现的"取消说成命中"用例转为通过 & 已完成 \\
把服务端与客户端的模式白名单做成同源常量 & 不再出现只改一侧的情况 & 1 小时 \\
删除或归档已停用的语音代码 & 仓库里不再有无法验证的能力声明 & 1 小时 \\
录一段完整演示 & 五条链路各有一段可播放记录 & 半天 \\
\bottomrule
\end{tabular}
\end{table}

最后一项需要说明：这台机器上没有可见桌面和音频，我生成的是逐步截图序列加文字说明（在 \code{output/demo/} 与 \code{docs/DEMO-WALKTHROUGH.md}），它不是录像。要让评审看到真实操作过程，需要人在正常桌面环境下录一次。

\section{复现}

\begin{verbatim}
npm start                      # 启动，默认 127.0.0.1:8765
npm test                       # 131 项自动测试
npm run build:knowledge        # 从引擎重新生成知识卡
npm run eval:retrieval         # 检索对照实验
npm run train:intervention     # 干预策略实验
\end{verbatim}

外部来源与它们各自的使用边界记在 \code{knowledge/sources.json}、\code{docs/DOMAIN-RESEARCH.md}。官方文档、公开转载与我的推断分开标注。
